\def\stat{l2-13}





\def\tit{ПОДХОДЫ К РАЗРАБОТКЕ ТЕХНОЛОГИИ  АВТОМАТИЗИРОВАННОГО КАРТОГРАФИЧЕСКОГО 


ПРЕДСТАВЛЕНИЯ ИНДИКАТОРОВ В СИСТЕМАХ МОНИТОРИНГА}








\def\titkol{Подходы к разработке %


%технологии  автоматизированного картографического 


представления индикаторов в системах мониторинга} % научной деятельности и~их динамики в~ИТСМ РАН} 





\def\autkol{Д.\,А.~Никишин}





\def\aut{Д.\,А.~Никишин$^1$}





\titel{\tit}{\aut}{\autkol}{\titkol}





%{\renewcommand{\thefootnote}{\fnsymbol{footnote}}\footnotetext[1]


%{Работа выполнена при поддержке РФФИ, грант 10-01-00480.}}





\renewcommand{\thefootnote}{\arabic{footnote}}


\footnotetext[1]{Институт проблем информатики 


Российской  академии наук, dmnikishin@mail.ru}      





      


      \Abst{Рассмотрены проблемы методологии и технологические 


решения для автоматизированного построения картограмм индикаторов 


научной деятельности на примере информационно-технологической 


системы мониторинга Российской академии наук (ИТСМ РАН). Новыми аспектами 


раз\-ра\-ба\-ты\-ва\-емой технологии являются представление наборов 


индикаторов для структурных подразделений РАН и отображение 


динамики индикаторов. Приведены образцы 


ин\-фор\-ма\-ци\-он\-но-кар\-то\-гра\-фи\-че\-ских документов с учетом структуры РАН, а также примеры 


отображения динамики индикаторов.}


     


     \KW{тематические карты; диаграммы; системы мониторинга}


     


       \vskip 12pt plus 9pt minus 6pt








      \thispagestyle{headings}





      


            \label{st\stat}


            


            \vspace*{-9pt}


     


     \section{Введение}


     


     Основные методические и технологические решения, используемые 


для картографической презентации данных ИТСМ, были рассмотрены в~[1]. 


Данная работа является дальнейшим развитием этого направления и 


затрагивает вопросы представления данных с учетом территориальной и 


структурной специфики уч\-реж\-де\-ний РАН, а также отображения динамики 


показателей и анализа тенденций их изменения во времени.


     


     Разрабатываемая в ИПИ РАН с 2003~г.\ 


     ИТСМ  РАН предназначена для обеспечения сбора, хранения, актуализации и 


анализа определенного круга данных, имеющих место в деятельности РАН. 


К~ним относятся как нормативные, так и фактографические данные об 


организациях, подразделениях и сотрудниках РАН, носящие различный 


характер (например, кадровые, юридические, финансовые и~т.\,п.). Наиболее 


значимой их составляющей являются данные, ха\-рак\-те\-ри\-зу\-ющие 


     на\-уч\-но-тех\-ни\-че\-скую деятельность РАН~--- отражающие научные 


результаты (такие как публикации, патенты, заявки и отчеты). Категории 


информационных ресурсов ИТСМ РАН подробно представлены в~[2].


     


     Структура базы данных (БД) ИТСМ обеспечивает возможность регистрации 


и учета этих данных с привязкой в пространстве и во времени. Наличие 


адресной привязки данных ИТСМ обеспечивает выполнение запросов по 


территориальному принципу в виде списков (таблиц) различных показателей 


ИТСМ. Обеспечиваемая современными БД возможность регистрации и 


накопления массива разновременных фактов позволяет проследить их 


динамику и выявить тенденции ее развития и сопровождающие ее процессы. 


Все это позволяет оперативно получать информацию как о территориальном 


распределении, так и об изменении во времени различных индикаторов 


научной деятельности в целом по стране, по отдельным  территориальным 


образованиям внутри страны (федеральным округам, субъектам Федерации), 


по городам, а также по различным региональным подразделениям 


организационной структуры РАН. 


     


     Региональная составляющая в структуре РАН достаточно значительна: 


так, доля бюджетных ассигнований региональным отделениям (без учета 


региональных научных центров центрального подчинения) приближается к 


40\% общего объема финансирования РАН. Статистика по объектам 


интеллектуальной собственности в организационной структуре РАН 


показывает следующее: тематические отделения РАН имеют 1723~объекта, 


региональные центры центрального подчинения~--- 182, а региональные 


отделения~--- 2278, т.\,е.\ более 50\%. Поэтому картографический аспект 


может быть полезен при иллюстрации распределений некоторых видов 


результатов по регионам.


     


     Как было показано в~[1], технологической основой построения карт, 


отоб\-ра\-жа\-ющих подобную информацию, является внедрение в ИТСМ 


специального картографического модуля, включающего в себя  специальное 


программные средства и дополнительные картографические данные, 


обеспечивающие наглядное отображение в картографическом пространстве 


сведений, содержащихся в БД ИТСМ.


     


     Направлением дальнейшего совершенствования этого модуля является 


расширение его функциональных возможностей путем организации его 


картографических объектов в виде схемы, соответствующей разветвленной 


организационно-тер\-ри\-то\-ри\-аль\-ной структуре РАН. Технологически это 


предполагает создание специализированной бланковой картографической 


основы, отображающей разными условными знаками пространственное 


положение подразделений различного уровня организационной структуры 


РАН (отделений, научных центров и отдельных институтов). Это должно 


позволить создавать в автоматизированном режиме различные карты 


динамики показателей для подразделений РАН различного уровня. 


     


     В качестве методической основы работы системы, адаптированной к 


построению карт динамики, применима традиционная методика 


формирования картографических отчетов, реализованная как 


функциональная задача в большинстве современных ГИС-технологий.


     


     В качестве основы картографического отчета используется специально 


разработанный шаблон~--- бланковая карта территории России, 


отображающая пространственное расположение структурных подразделений 


РАН, включая местоположение элементов структуры самых низших 


порядков (институты, филиалы и~пр.), а в роли исходных входящих потоков 


данных для их заполнения~--- постоянно актуализируемые данные, 


поступающие по территориальным запросам из БД ИТСМ. 


     


     Проблемами реализации подобной функциональности являются:


     \begin{enumerate}[1.]


     \item Сложная, разветвленная организационная структура РАН, 


предполагающая наличие нескольких классификаций.


\item Значительный разброс значений индикаторов научной деятельности  


для различных подразделений и учреждений РАН, обусловленный их 


территориальной спецификой.


     \item Неравномерность в степени заполнения ИТСМ РАН данными (в 


частности, обеспеченности данными различных территориальных и 


организационно-струк\-тур\-ных образований).


     \item В перспективе~--- обеспечение возможности использования для 


анализа данных и вычисления индикаторов и оценок эффективности 


подразделений РАН дополнительных данных об индивидуальных индексах 


цитирования различных категорий сотрудников.


     \end{enumerate}


     


     Для практической реализации разработки требуется разработать 


универсальную модель организационной структуры РАН, допускающей 


несколько параллельных классификаций, и соответствующую ей систему 


шаблонов карт для картографического отображения результатов анализа, 


проводимого по различным подразделениям РАН.


     


     \section{Методологические и~технологические приемы создания 


по~данным~информационно-технологической системы мониторинга~картографических  документов 


с~учетом структуры~Российской~академии~наук}


     


     Специфика подготовки данных для автоматического формирования 


динамических карт определяется организационной структурой учреждений 


РАН~[3], характеризующейся наличием нескольких классификаций:


     \begin{itemize}


     \item по отделениям РАН по направлениям наук;


     \item по региональным отделениям РАН;


     \item по региональным научным центрам.


     \end{itemize}


     


     Схема подчинения основных элементов организационной структуры 


РАН приведена на рис.~1. В~схему не вошли руководящие и специальные 


органы РАН (президиумы, комитеты, научные советы и~пр.), 


присутствующие в каждой из ветвей структуры.


     


     Наличие связей типа $1:\infty$ (<<один ко многим>>) между разными 


составляющими организационной структуры (отделениями РАН, научными 


центрами отделений, региональными центрами центрального подчинения 


и~пр.) обеспечивает однозначный переход к спискам подчиненных 


элементов структуры~--- институтов (учреждений) и их филиалов, в свою 


очередь предоставляющих данные в БД ИТСМ. 


\begin{figure} %fig1


       \vspace*{1pt}


\begin{center}


\mbox{%


\epsfxsize=129.28mm


\epsfbox{ipi-21(2)-13-1.eps}


}


\end{center}


\vspace*{-9pt}


\Caption{Схема подчинения элементов организационной структуры РАН}


\end{figure}


Поскольку все показатели 


(данные) в ИТСМ привязаны по адресу предоставившей их организации, 


получение консолидированных данных по каждому показанному на схеме 


структурному подразделению РАН не представляет особых трудностей. Это  


осуществляется построением запросов, предоставляющих результаты 


суммирования показателей, полученных от совокупности подчиненных 


организаций той или иной структуры РАН.


     


\begin{figure} %fig2


       \vspace*{1pt}


\begin{center}


\mbox{%


\epsfxsize=125.855mm


\epsfbox{ipi-21(2)-13-2.eps}


}


\end{center}


\vspace*{-9pt}


\Caption{Образец карты-основы России с  локализацией местоположения 


структурных подразделений РАН}


\end{figure}


     


     


     Для картографической визуализации данных подобных запросов 


необходимо на специальной цифровой карте-бланковке со схематичным 


изображением территории России пространственно локализовать 


местоположение всех организаций, входящих в структуру РАН. Как правило, 


это осуществляется привязкой крупных подразделений РАН (отделений, 


региональных центров и~пр.) к крупным городам, в которых они базируются. 


На представленной карте-основе, фрагмент которой представлен на рис.~2, 


разными условными знаками (тематическими картографическими объектами) 


представлены основные структурные подразделения РАН, включая:


     \begin{itemize}


     \item три региональных отделения РАН (Уральское, Сибирское и 


Дальневосточное); 


     \item входящие в них научные центры отделений;


     \item региональные научные центры центрального подчинения.


     \end{itemize}


     


     Обязательным условием является обязательное приписывание в 


семантике картографических объектов, соответствующих различным 


элементам структуры РАН, названий, полностью соответствующих 


названиям в БД ИТСМ. Именно по названию организации\,--\,поставщика 


данных в ИТСМ осуществляется связь показателей с картой.


     





     При необходимости (например, для осуществления картографической 


презентации по направлениям наук) особыми знаками можно отобразить 


ареалы сосредоточения профильных институтов, подчиняющихся 


отделениям по различным направлениям наук (как правило, отделения по 


направлениям наук сосредоточены в крупных городах~--- Москве,  Санкт-Петербурге и некоторых~других).


     


     Картографическая презентация показателей ИТСМ может 


производиться также для объектов самых нижних уровней организационной 


структуры РАН~--- отдельных институтов (учреждений) и/или их филиалов. 


Для этой цели в отдельном слое бланковой карты должно быть осуществлено 


позиционирование каждого из них в соответствии с его принадлежностью к 


городу его расположения, прописанному в адресной привязке.


     


     Наличие трех ветвей в организационной структуре позволяет строить 


различные карты, отображающие величину различных показателей по годам 


и динамику их изменения во времени. Построение карт может 


осуществляться как отдельно, для различных подразделений РАН, так и 


консолидированно, для всей совокупности выделяемых пользователем 


картографических объектов, соответствующих различным элементам 


структуры РАН, представленным на карте.


     


     С этой целью из данных единой бланковой карты России средствами 


ГИС (геоинформационной системы) в полуавтоматическом режиме формируются специализированные 


карты, используемые для отображения результатов проводимого анализа. 


Пример подобной карты, сформированной для отображения показателей БД 


ИТСМ применительно к региональным центрам центрального подчинения, 


представлен на рис.~3. Аналогичные карты могут быть сформированы на 


региональные отделения РАН, хотя вся работа может осуществляться и по 


исходной карте России, а для представления результатов осуществлять 


вырезание нужного ее фрагмента.


     





     


     Аналогично могут быть сформированы картографические основы, 


необходимые для построения динамических карт для отделений РАН по 


направлениям наук. Подобные основы для каждого отделения должны 


отображать места локализации институтов, подчиняющихся данному 


отделению. Главным образом подобные институты локализованы в крупных 


городах. 


     


     Прописанное в структуре РАН многоуровневое соподчинение объектов 


позволяет осуществлять картографическую  презентацию показателей ИТСМ 


в соответствии с географическим положением как каждого из структурных 


подразделений РАН (например, презентация по пространственной  


принадлежности региональных отделений РАН), так и  по территориальной 


принадлежности отдельных учреждений РАН (научных центров, институтов 


и учреждений-филиалов). Аналогично  можно проанализировать индикаторы 


научной деятельности для коллективов журналов, научных советов, 


комитетов и~пр., подчиняющихся тому или иному 


     организационно-структурному подразделению РАН либо входящим в 


определенный субъект Российской Федерации.





     \begin{figure} %fig3


            \vspace*{1pt}


\begin{center}


\mbox{%


\epsfxsize=118.705mm


\epsfbox{ipi-21(2)-13-3.eps}


}


\end{center}


\vspace*{-9pt}


     \Caption{Картографическая основа для презентации данных ИТСМ в 


части региональных центров РАН}


     \end{figure}


     


     Все числовые данные для построения карт  состояния и динамики 


показателей на\-уч\-но-интеллектуальной деятельности поставляются из БД 


ИТСМ. Как было показано в~[1], при создании карт показателей ИТСМ на 


различные регионы России это осуществлялось посредством запросов, в 


которых название региона,  сформированное на основе адресной привязки 


показателей ИТСМ, являлось связующим полем с картой. В~нашем случае 


построения карт применительно к различным подразделениям 


организационной структуры РАН основным связующим полем в запросах 


является название института (учреждения) РАН, поставившего данные в 


ИТСМ.  Фрагмент результата типового запроса из БД ИТСМ представлен на рис.~4.


     


\begin{figure} %fig4


       \vspace*{1pt}


\begin{center}


\mbox{%


\epsfxsize=120mm


\epsfbox{ipi-21(2)-13-4.eps}


}


\end{center}


\vspace*{-9pt}


\Caption{Фрагмент результата запроса из БД ИТСМ}


\end{figure}





     Общим требованием к запросам является:


     \begin{itemize}


     \item наличие связующего поля, по значению которого осуществляется 


связь БД с территориальными образованиями (объектами) карты;


     \item наличие одного или нескольких полей  со значениями показателей, 


подлежащих отображению на тематической карте/\!карте-диаграмме. 


     \end{itemize}


     


     Для частного случая разработки это будут, соответственно, поле 


NAME\_RUS с именами подразделений (отделений РАН и научных центров) 


и несколько полей значений, содержащие, например, данные о количестве 


заявок за разные годы (поля DCL\_2003, DCL\_2004 и~т.\,д.), подлежащие 


отображению на тематических картах. Поля запросов могут содержать также 


данные об изменении показателей за разные годы, что позволяет создавать 


карты-диаграммы, отражающие тенденции изменения показателей во 


времени.


     


     \section{Создание картографической презентации информационно-технологической системы 


     мониторинга~с~использованием ГИС-приложений}


     


     Основные методические и технологические приемы, используемые для 


картографической презентации данных ИТСМ, были рассмотрены в~[1] и в 


применении к созданию карт динамики индикаторов сводятся к созданию 


дополнительной картографической основы, отображающей 


пространственное положение объектов организационной структуры РАН.


     


     Автоматизированное построение электронной тематической карты 


включает в себя два основных этапа:


     \begin{enumerate}[1.]


     \item Выборку из БД ИТСМ требуемых тематических показателей (в 


виде массива, содержащего значения одной или нескольких характеристик 


для каждого из пространственных объектов). Данный этап может быть 


реализован стандартными методами запросов из БД (таких как SQL-Server, 


Access и~т.\,п.). Для обеспечения корректного взаимодействия с БД должно 


быть предусмот\-ре\-но согласование значений ключевых полей БД (названий 


пространственных объектов) с соответствующими именами объектов 


картографической основы.


     \item Формирование тематической карты~--- отображение 


тематических показателей на фоне базовой картографической основы 


(используемой как подложка). Для реализации данного этапа требуются 


средства, позволяющие автоматизированно наносить и редактировать линии, 


полигоны, точечные знаки, подписи на растровую или векторную 


картографическую основу. В~за\-ви\-си\-мости от используемого инструментария 


и условий дальнейшего использования (внедрения в документы-отчеты) 


создаваемая тематическая карта должна представляться одном из следующих 


видах:


     \begin{itemize}


     \item в стандартных картографических форматах (SXF/RSW, MIF/MID, 


SHP  и~т.\,п.);


     \item в векторных форматах общего назначения (DXF, SVG);


     \item в виде растрового изображения карты (BMP, TIFF, JPEG, PNG, 


GIF и~т.\,п.).


     \end{itemize}


     


     Для отображения тематической карты в браузере или для ее внедрения 


в документ наиболее оптимальным представляется последний вариант.


     \end{enumerate}





          \begin{figure}[b] %fig5


            \vspace*{-9pt}


\begin{center}


\mbox{%


\epsfxsize=123.36mm


\epsfbox{ipi-21(2)-13-5.eps}


}


\end{center}


\vspace*{-9pt}


     \Caption{Картограмма с отображением количества заявок, поданных 


от региональных центров в 2007~г. }


     \end{figure}


     


     Исходные данные для формируемых карт запрашиваются 


непосредственно из БД ИТСМ. В~данном случае производится выделение 


всех объектов, отображающих необходимые для анализа подразделения РАН. 


Возможно построение карт по всему множеству организационных структур 


РАН, включая отделения РАН, региональные научные центры и научные 


центры отделений. В~результате каждый выбранный для построения объект 


организационной структуры, для которого нашлась соответствующая запись 


в БД ИТСМ, отображается цифровым показателем из БД. 


     \begin{figure}[b] %fig6


            \vspace*{-9pt}


\begin{center}


\mbox{%


\epsfxsize=124.813mm


\epsfbox{ipi-21(2)-13-6.eps}


}


\end{center}


\vspace*{-9pt}


     \Caption{Секторная картодиаграмма с отображением количества 


заявок, поданных от региональных центров за несколько лет 


     (2006--2008~гг.)}


     \end{figure}


     Так, например, на 


рис.~5 каждый региональный научный центр РАН отображается 


соответствующим числовым показателем количества заявок, поданных в 


2007~г. Плотность цветовой заливки цифровых показателей отвечает 


определенной градации, определяемой величиной показателя согласно 


легенде картодиаграммы.


     





     


     Аналогичным образом могут быть построены картодиаграммы, 


отобра\-жа\-ющие для различных объектов структуры РАН прирост показателей 


(в сравнении с прежними годами). Это позволяет выделить на карте разными 


цветами благополучные и проблемные научные центры, отобразив цветом 


разные тенденции изменения показателя во времени (рис.~6--8).


     





     


     \begin{figure} %fig7


            \vspace*{1pt}


\begin{center}


\mbox{%


\epsfxsize=125.35mm


\epsfbox{ipi-21(2)-13-7.eps}


}


\end{center}


\vspace*{-9pt}


     \Caption{Прирост числа заявок в 2008~г.\ по региональным научным 


центрам (относительно предыдущего года)}


     \end{figure}





\begin{figure} %fig8


       \vspace*{1pt}


\begin{center}


\mbox{%


\epsfxsize=125.811mm


\epsfbox{ipi-21(2)-13-8.eps}


}


\end{center}


\vspace*{-9pt}


\Caption{Прирост числа проектов в 2008~г.\ по научным центрам 


Уральского и Сибирского отделений РАН относительно предыдущего года}


\end{figure}


     


     Идентичность названий полей таблиц БД и ключей семантики объектов 


карты обеспечивает возможность идентификации объекта карты и записи 


БД (связь таб\-ли\-цы и карты), что позволяет автоматизировать 


процессы привязки БД к карте и обновления объектов карты. Для 


обеспечения этой возможности необходимо, чтобы названия, например, 


регионов РФ в семантическом описании объектов исходной 


картографической основы и в таблицах тематических показателей были 


абсолютно идентичны.


     


     В качестве источника информации о пространственном расположении 


анализируемых объектов ИТСМ~--- организаций~--- используется их почтовый 


адрес. Эта информация содержится в БД <<Рубрикатора стран, регионов, 


городов>>, представляющего собой иерархическую структуру с 


наименованиями федеральных округов, регионов и городов РФ, а также 


регионов и городов других стран. Структура рубрикатора может дополняться 


и уточняться по мере пополнения и уточнения Адресной~БД. 


     


     В качестве нормативных документов для наименований субъектов в 


классификаторе была использована статья~65 Конституции РФ, а также 


ГОСТ 7.67-2003 <<Коды для представления названий стран>> и аутентичный 


ему ISO~3166-88, устанавливающие  буквенные и цифровые обозначения 


названий стран в кодированной форме, единые для различных систем 


обработки информации, ее хранения и обмена. В~случае расхождений 


названий в БД и на карте за истину должно приниматься название в БД 


(исправления производятся в названии региона на карте).


     


     Перед непосредственным созданием карты следует проанализировать 


структуры таблиц показателей, которые будут использоваться для создания 


карты и определить:


     \begin{itemize}


     \item перечень тематических показателей, подлежащих визуализации 


на карте (названия полей и таблиц этих показателей в~БД);


     \item способ визуализации объектов тематической карты, в том числе 


цвет, раз\-мер, точ\-ность численных показателей (количество знаков после 


запятой), взаимный порядок расположения, а также вид картодиаграммы.


     \end{itemize}


     


     На основе этого выполняется настройка отображения объектов 


тематической карты (условных знаков, которые будет использованы при 


создании карты) путем изменения соответствующих атрибутивных описаний. 


Такая настройка обычно выполняется в интерактивном режиме, однако в 


условиях автоматизированного создания тематических карт необходимо 


обеспечение автоматической настройки изображения объектов (их 


атрибутивных описаний) в зависимости от вида создаваемой карты.


     


     Все объекты, отображаемые на карте, должны быть логически 


распределены на слои отображения, определяющие приоритетность вывода 


на экран (чем выше приоритет, тем слой будет выводиться позже, 


соответственно, в меньшей степени буден перекрываться графикой слоев с 


более высоким приоритетом, выводимых позже).


     


     Оперативное обновление (ведение) существующей тематической карты 


подразумевает автоматическое изменение изображений объектов 


тематической карты при изменении в БД соответствующих им значений. 


После изменения семантических характеристик объекта его изображение на 


карте должно изменяться автоматически в соответствии с новым значением. 


Однако, поскольку в условиях ИТСМ использование тематических карт 


носит эпизодический характер, в виде статических документов, 


необходимости в обеспечении оперативного обновления нет.


     


     Варианты реализации способов автоматизированного построения 


тематических карт:


     \begin{enumerate}[1.]


     \item Использование инструментов тематического картографирования, 


имеющихся в составе ряда ГИС, включающих средства построения запросов 


к внешним БД и отображения их результатов на картографической основе. 


Примеры такого подхода~--- инструменты тематического картографирования 


в ГИС MapPoint и <<Панорама>>~[1]. Необходимость наличия 


установленного программного обеспечения ГИС, знания и навыки работы с ГИС сильно 


ограничивают возможности их использования для целей ИТСМ.


     \item Противоположным вариантом является создание собственного 


модуля тематического картографирования исключительно силами 


разработчика. Однако, если работа по извлечению данных из БД особых 


сложностей не вызывает, то самостоятельная разработка модуля для создания 


карты является весьма непростой задачей, требующей специфических знаний 


и навыков работы с растровой и векторной графикой.


     \item Наиболее приемлемым (по функциональности и сложности 


реализации) является вариант с внедрением в разрабатываемое приложение 


(в данном случае в ИТСМ) готовых средств работы с картами (и даже 


средств по\-стро\-ения тематических карт). Такие средства предоставляются в 


составе того или иного <<инструментария разработчика приложений>>. 


В~качестве примера можно привести:


     \begin{itemize}


     \item <<компонент географического отображения данных>> из пакета 


Crystal Reports, предназначенный для использования в офисных 


приложениях~[4]. Он использует идеологию ГИС <<MapInfo>>, 


использованную и в ГИС MS MapPoint;


     \item <<Инструментарий разработчика ГИС-приложений GIS 


ToolKit>>~[5], в частности GIS WebToolKit~--- инструментарий  для 


разработки Web-сай\-тов на платформе ASP.NET в среде .NET Framework 3.5  


средствами Microsoft Visual Studio 2008, язык программирования  C\#. GIS 


WebServer отображает на карте, снимке данные об объектах БД, 


имеющих территориальную привязку. В~приложении возможно 


использование различных БД: MS SQL Server, Oracle, MS Access и 


других. Работа с картами и БД в GIS WebServer выполняется в 


браузере любого типа: MS Internet Explorer, Mozilla (FireFox), Opera 


и~другие. Изображение карты создается динамически и выводится в 


     Web-браузер в виде png-рисунка (рис.~9).


     \end{itemize}


     \end{enumerate}


          Их практическому использованию в ИТСМ в настоящий момент 


препятствуют следующие моменты:


     \begin{itemize}


     \item в имеющемся варианте Crystal Reports <<компонент 


географического отоб\-ра\-же\-ния данных>> недоступен, поскольку требует 


дополнительного лицензирования; кроме того, все карты, 


предоставляемые программой Crystal Reports, являются англоязычными (но 


при этом есть возможность установки соответствий русскоязычных названий 


городов англоязычным названиям); карты разбиты по континентальному 


принципу, поэтому нет отдельной карты на территорию Российской 


Федерации. 


     \item <<Инструментарий разработчика ГИС-приложений GIS ToolKit>> 


пред\-став\-ля\-ет собой набор компонентов ГИС <<Панорама>>. Лицензируется 


либо его использование без ограничения распространения (159\,000~руб.), 


либо распространение использующих его приложений с электронным 


ключом (5100~руб.). Он является средством, специализированным под 


решение специфических задач ГИС, более трудоемким, проблематичным 


в плане его интеграции с ИТСМ, поэтому его использование для целей 


ИТСМ представляется менее целесообразным.


      \end{itemize}





\begin{figure} %fig9


\vspace*{1pt}


\begin{center}


\mbox{%


\epsfxsize=125mm


\epsfbox{ipi-21(2)-13-9.eps}


}


\end{center}


\vspace*{-10pt}


\Caption{ Выполнение демонстрационного примера GIS WebToolKit~--- 


Web-сервер, предназначенный для интерактивного взаимодействия с 


клиентами}


\vspace*{-3pt}


\end{figure}





\vspace*{-6pt}


     


     \section{Заключение}


     


     \vspace*{-1pt}


     


     Затронутые в данной работе вопросы представления данных с учетом 


территориальной и структурной специфики учреждений РАН, а также вопрос 


отображения динамики показателей и анализа тенденций их изменения во 


времени становятся актуальны для определенных задач отображения 


индикаторов, имеющих место в деятельности РАН. Это должно позволить 


создавать в автоматизированном режиме различные карты динамики 


показателей для подразделений РАН различного уровня, выявить тенденции 


их изменения и сопутствующие этому процессы.


     


     Основными проблемами реализации подобной функциональности 


являются:


     \begin{itemize} %[1.]


     \item сложная, разветвленная организационная структура РАН;


     \item значительный разброс значений индикаторов научной 


деятельности  для различных подразделений и учреждений РАН, 


обусловленный их территориальной спецификой;


     \item неравномерность в степени заполнения ИТСМ РАН данными (в 


частности, обеспеченности данными различных территориальных и 


организационно-структурных образований).


     \end{itemize}


     


     Для практической реализации разработки требуется разработать 


универсальную модель организационной структуры РАН, допускающей 


несколько параллельных классификаций, и соответствующую ей систему 


шаблонов карт для картографического отображения результатов анализа, 


проводимого по различным подразделениям РАН.


     


     {\small\frenchspacing


{%\baselineskip=10.8pt


\addcontentsline{toc}{section}{Литература}


\begin{thebibliography}{9}





      \bibitem{1-2-13}


      \Au{Никишин~Д.\,А.}


      Методические и технологические решения для картографического 


представления индикаторов научной деятельности в 


      информационно-технологической системе мониторинга РАН~/\!\!/ 


Системы и средства информатики.~--- М.: ИПИ РАН, 2010.   Вып.~20. №\,2.


С.~100--112.


      


      \bibitem{2-2-13}


      \Au{Лунева~Н.\,В.}


      Категории информационных ресурсов ИТСМ РАН~/\!\!/ Системы и 


средства информатики.~--- М.: ИПИ РАН, 2010.  Вып.~20. №\,2. С.~88--99.


      


      \bibitem{3-2-13}


      Официальный сайт Российской академии наук. Структура РАН. {\sf 


http://www.ras.ru/}.


      


      \bibitem{4-2-13}


      \Au{Маклаков С.\,В., Матвеев Д.\,В.}


      Анализ данных. Генератор отчетов Crystal Reports.~--- СПб.: 


      БХВ-Петербург, 2003.  496~с.


       


       \label{end\stat}


       


      \bibitem{5-2-13}


      GIS Web Tool Kit~--- инструментарий для разработки 


      Web-приложений. {\sf 


http:// www.gisinfo.ru/products/giswebtoolkit.htm}.


\end{thebibliography}


}


}    


